\section{Limitations}
\label{sec:conclusion_limitations}

Although the research presented in this thesis has produced conclusive results, it remains subject to limitations.
This section discusses the five main limitations identified at this stage.

\gls{car} describes how the \gls{ls} of an agent is organized, but not what each cluster means.
The metrics quantify the consistency and abstraction of the clusters, but they say nothing about their content.
In this manuscript, the meaning of each cluster was assigned through manual inspection of the corresponding observations.
This step is time-consuming, subjective, and difficult to scale to a large number of clusters or environments.

The clusters identified by \gls{car} are not yet directly related to the decisions of the agent.
We observe that the agent distinguishes certain situations, but not how these distinctions lead to a given action.
As a result, \gls{car} cannot yet explain a specific failure or reveal a bias in the behavior of the agent.

The experiments were conducted on six Atari environments with relatively small networks.
These environments are diverse, but they remain limited in size and complexity.
Whether \gls{car} scales to more complex environments and larger networks has not yet been established.
In addition, only one clustering method was evaluated, based on $k$-means with Euclidean distance.
Other distances and other clustering methods could reveal different structures in the \gls{ls}.
Since the pipeline of \gls{car} is modular, these alternatives can be tested and compared using the proposed metrics.

The proposed metrics have only been applied to clustering.
The other families of methods reviewed in Chapter~\ref{chap:latent_space} have not yet been evaluated within this framework.
This is particularly the case for methods based on the \SuperpositionHypothesisName, such as \glspl{sae}.
These methods are widely used, but their assumptions remain debated.
Our metrics could be used to test these assumptions empirically and quantitatively, and to compare these methods with \gls{car} on the same tasks.
